\subsection{The quotient topos of a foliation: 119/145--146}

We use the term \emph{\'etendue} for a Grothendieck topos $\mathcal T$
with an epimorphism $V\to1$ for which $\mathcal T/V$ is localic.
This is the topos interpretation of the manuscript's topological
multiplicity. A discrete group $G$ gives the \emph{\'etendue} $BG=G\text{-}\mathrm{Set}$:
the regular $G$-set is such a $V$ and $BG/V\simeq\mathrm{Set}$.
The corresponding stack is the quotient stack $[*/G]$;
we do not identify a stack with its coarse orbit space.

A local equivalence relation on a space or locale $X$ is a germwise
compatible atlas of equivalence relations $R_i\subset U_i\times U_i$.
The two projections are required to be open. We require arbitrarily
fine atlases with connected, locally connected projections. An invariant
sheaf has transport for these relations, locally on $X$, with the unit
and cocycle conditions. Connectedness makes transport unique if it
exists. Thus $\operatorname{Sh}(X,r)$ is a full subcategory of
$\operatorname{Sh}(X)$. In this generality it need not be a Grothendieck
topos; local simple connectivity is a sufficient additional hypothesis
\cite[Example 2.3; Theorem 5.3]{KockMoerdijk1992}.

\begin{saresult}{A.701}{Refuted}{Connected fibres do not determine the monodromy quotient}
Interpret $\mathcal C'$ in 119/146 as triples $(X,\mathcal T,p)$,
where $\mathcal T$ is an \emph{\'etendue} and
$p:\operatorname{Sh}(X)\to\mathcal T$ is open, connected and locally
connected. Sending such a triple to its local kernel does not give an
equivalence with local equivalence relations. This remains false for
the circle with its single-leaf relation, with every fibre a connected
one-dimensional manifold. The manuscript's third condition is indeed
``et connexes'', as checked on the original scan, not simple connectivity.
\end{saresult}
\begin{proof}
Let $X=S^1$ and let $r$ have the single chart $X\times X$; on every
small arc it is the indiscrete relation. For $n\ge2$ let
$G_n=\Z/n\Z$ and let $P_n=S^1\to S^1$ be the principal $G_n$-cover
$z\mapsto z^n$. Its classifying geometric morphism
\[
 p_n:\operatorname{Sh}(S^1)\longrightarrow BG_n
\]
has inverse image sending a $G_n$-set $A$ to the sheaf of sections of
$P_n\times^{G_n}A$. Pullback along the epimorphic \emph{\'etale} atlas
$\mathrm{Set}\to BG_n$ gives $\operatorname{Sh}(P_n)\to\mathrm{Set}$.
This map is open and locally connected since $P_n$ is a locally
connected space, and connected since $P_n$ is connected. These
properties descend along this atlas. More concretely, full faithfulness
of $p_n^*$ follows because a map of the associated local systems is a
map of their fibres commuting with the monodromy of $\pi_1(S^1)=\Z$,
and the homomorphism $\Z\to G_n$ is onto.

We may also take $G_1=1$, $P_1=S^1$, and the structural map $p_1$ to
$\mathrm{Set}$. On an arc trivializing $P_n$, the map to the atlas is
the structural map of that arc. Its kernel relation is the indiscrete
relation on the arc. Hence $\operatorname{Lker}(p_n)=r$ for every $n$,
including $n=1$. Nevertheless $BG_2$ is not equivalent to
$\mathrm{Set}$: its regular two-element $G_2$-set is connected and has
two automorphisms, whereas a noninitial connected object of
$\mathrm{Set}$ is a singleton and has only the identity automorphism.
Thus the local-kernel construction is not injective on equivalence
classes of objects above this fixed $X$.

The universal cover $\R\to S^1$ similarly gives
$p_\infty:\operatorname{Sh}(S^1)\to B\Z$. Its inverse image has as
essential image all locally constant sheaves. To check this directly,
pull a local system back to $\R$; path lifting trivializes it, and deck
translation gives the action of $\Z$ on a fibre. Conversely the
associated bundle of any $\Z$-set is a local system, and these two
constructions are inverse on objects and morphisms. Invariant sheaves
for $r$ are exactly these local systems: on each small arc, transport
for the indiscrete relation is precisely a constant trivialization.
Therefore $\operatorname{Sh}(X,r)\simeq B\Z$, the monodromy quotient.
The coarse, or codimension-zero holonomy, quotient is a point.

We can also classify all the discrete-group examples just used.
A principal discrete $G$-cover $P\to S^1$ with connected $P$ has
monodromy given by a surjection $\Z\to G$: transitivity on the regular
$G$-fibre is equivalent to surjectivity. Thus $G$ is cyclic, finite or
infinite, and the above construction gives all such presentations up
to an automorphism of $G$. This is a classification in the explicitly
stated discrete-group class, not a classification of all \emph{\'etendues}.
\end{proof}

\begin{saresult}{A.702}{Proved}{The correctly based Hom groupoids are essentially discrete}
For two triples $(X,\mathcal T,p)$ and $(Y,\mathcal U,q)$ as above,
define a morphism to be $(f,g,\alpha)$, where $f:X\to Y$ is continuous,
$g:\mathcal T\to\mathcal U$ is a geometric morphism, and
$\alpha:q\operatorname{Sh}(f)\Rightarrow gp$ is an invertible
$2$-cell. A $2$-cell between such morphisms with the same $f$ is a
$2$-cell $\theta:g\Rightarrow g'$ compatible with their displayed
triangles. There are no $2$-cells between different $f$'s.
If $p$ is connected, each such Hom category is equivalent to a discrete
category. It need not be literally discrete before taking isomorphism
classes. Omitting the triangle changes this conclusion.
\end{saresult}
\begin{proof}
The compatibility equation determines $\theta p$ uniquely from
$\alpha,\alpha'$. On inverse image functors this determines
$p^*\theta^*$, where $p^*$ is fully faithful. Faithfulness implies
that there is at most one possible component $\theta_A^*$ for each
object $A$ of $\mathcal U$, hence at most one $2$-cell. Since its
image under $p^*$ is invertible, full faithfulness also implies that
each component is invertible. Thus the Hom category is a groupoid
with at most one arrow between any two objects; choosing its set of
isomorphism classes gives an equivalent discrete category. Every
automorphism is the identity. If we omit the triangle, the identity
geometric morphism of $BG$ has automorphisms given by central elements
of $G$: on a $G$-set, act by that central element. Taking $G=\Z/2\Z$
gives a nonidentity automorphism. This verifies the distinction.
\end{proof}

\begin{saresult}{A.703}{Proved under stated hypotheses}{Exact recognition and a sufficient simple-connectivity repair}
Let $r$ be a locally connected local equivalence relation on $X$ for
which $\mathcal M_r=\operatorname{Sh}(X,r)$ is a Grothendieck topos
and the inclusion is the inverse image of
$\pi_r:\operatorname{Sh}(X)\to\mathcal M_r$.
For a connected $p:\operatorname{Sh}(X)\to\mathcal T$, there is an
equivalence $\mathcal T\simeq\mathcal M_r$ over
$\operatorname{Sh}(X)$ if and only if
\[
 \operatorname{EssIm}(p^*)=\operatorname{Sh}(X,r).
 \tag{A.7.1}
\]
This is the exact additional descent condition; connected fibres alone
do not imply it. The following geometric condition suffices. Suppose
$\mathcal T$ is an \emph{\'etendue}, $p$ is connected and locally
connected, $r=\operatorname{Lker}(p)$, and there is an epimorphic
\emph{\'etale} localic atlas $\operatorname{Sh}(A)\to\mathcal T$
whose pullback is $\operatorname{Sh}(B)\to\operatorname{Sh}(X)$,
such that the resulting map of locales $b:B\to A$ is simply connected
in the following explicit sense: $b^*$ is fully faithful, and any
sheaf on $B$ locally isomorphic to pullbacks by $b^*$ is globally such
a pullback. Then (A.7.1) holds. In particular $\mathcal M_r$ is an
\emph{\'etendue} in this situation.
\end{saresult}
\begin{proof}
For the first assertion both inverse image functors are full and
faithful. Equality of their essential images gives an equivalence
between their domains, intertwining the two functors. It preserves
finite limits and all colimits because the two inverse image functors
do and reflect isomorphisms; equivalently its quasi-inverse is an
equivalence of topoi. Necessity is immediate by applying inverse
image to an equivalence over $\operatorname{Sh}(X)$.

We check the sufficient condition by descent, explaining the comparison
map rather than identifying topoi from their points. Let $e:B\to X$
be the pulled-back atlas. The square of topoi with vertices
$\operatorname{Sh}(B),\operatorname{Sh}(A),\operatorname{Sh}(X),
\mathcal T$ is also a pushout. Indeed the two vertical atlas maps
are effective descent morphisms; their nerves compute the two lower
topoi, and the maps between corresponding nerve terms are pullbacks
of the connected, locally connected map $p$. Descent of arrows along
these faithful inverse images identifies the pushout with $\mathcal T$.
This is exactly \cite[Lemma 7.2, using Lemma 3.5]{KockMoerdijk1992};
its proof requires the stated connectedness, not contractibility of
an arbitrary pointwise fibre.

The simple-connectivity assumption gives
$\operatorname{Sh}(A)\simeq
\operatorname{Sh}(B,\operatorname{Lker}(b))$ by
\cite[Lemma 5.1]{KockMoerdijk1992}. Moreover
$\operatorname{Lker}(b)=e^*r$ in the sense of pulled-back local
relations, by that paper's Lemma 6.4. The pushout is the category of
triples $(F,D,\gamma)$ with $F\in\operatorname{Sh}(X)$,
$D\in\operatorname{Sh}(A)$ and
$\gamma:e^{-1}F\xrightarrow{\sim}b^{-1}D$.
Full faithfulness of $b^{-1}$ says that this category is the full
subcategory of sheaves $F$ for which $e^{-1}F$ is an $e^*r$-sheaf.
The existence of local transport descends along the surjective local
homeomorphism $e$; its uniqueness ensures compatibility on overlaps.
It is therefore equivalent to $F$ being an $r$-sheaf. This proves
(A.7.1) and identifies the comparison functor. The atlas exhibits the
resulting topos as the given \emph{\'etendue} $\mathcal T$.
\end{proof}

\begin{saresult}{A.704}{Proved under stated hypotheses}{A precise categorical and relative version}
Restrict objects on the quotient side to presentations $(X,r,\mathcal T,p)$
remembering the recognized relation and satisfying (A.7.1), and on the relation side to relations whose invariant-sheaf
category is an \emph{\'etendue} with its connected quotient map.
Define a morphism $f:(X,r)\to(Y,s)$ to be a continuous map for which
$f^{-1}$ takes $s$-invariant sheaves to $r$-invariant sheaves. Then
the quotient construction is an equivalence from this category to
the $2$-category of such presentations after replacing its Hom
groupoids by their isomorphism classes. These are explicitly
\emph{sheaf-detected} morphisms; a theorem for a different notion of
foliated map is not being asserted without a comparison.
The relation decoration can be omitted in the subclass with
$r=\operatorname{Lker}(p)$. On the relation side it suffices to require
locally simply connected effective charts: each chart relation is the
kernel pair of its quotient map. Then the local kernel of the canonical
quotient is the original relation, by \cite[Remark 6.5]{KockMoerdijk1992}.
Without effectivity one must not assert that equality merely from
equality of the invariant-sheaf categories.

The same recognition criterion and equivalence hold relative to a
Grothendieck topos $\mathcal E$, for internal locales, internal local
relations and relative topoi, with the foregoing descent and
simple-connectivity hypotheses imposed internally and stably under
base change. The unrestricted suggestion ``[ou topos ?]'' is false,
since its instance over $\mathrm{Set}$ is A.701.
\end{saresult}
\begin{proof}
For an allowed $f$, restriction of $f^{-1}$ is a finite-limit and
colimit preserving functor between the invariant-sheaf topoi.
It has the right adjoint
$\pi_{s*}f_*\pi_r^*:\mathcal M_r\to\mathcal M_s$.
The adjunction follows directly from the three adjunctions and full
faithfulness of the two inclusions. This constructs $g$ and
the triangle $\alpha$. Conversely such a triangle forces $f^{-1}$
to preserve the two essential images. Connectedness and A.702 give
uniqueness up to a unique compatible $2$-cell. A.703 supplies essential
surjectivity on the stated objects. Composition is restriction of
inverse images and its associativity constraint, so it gives the
asserted equivalence.

For the relative claim, use sheaves internal to $\mathcal E$ in
every construction. The atlas is an internal effective epimorphism;
descent data are still objects with an isomorphism satisfying a unit
and cocycle equation. Unique transport is proved by full faithfulness
of the connected inverse image, exactly as above, and glues along
the atlas. Hence the triple description in A.703 and its comparison
functor remain valid. Finite limits, colimits, inverse image
adjunctions and full faithfulness have the same relative meanings.
These arguments prove the recognition criterion and the equivalence
under the stated stable internal hypotheses; they do not infer
simple connectivity from connected geometric fibres. The circle
counterexample over $\mathrm{Set}$ proves falsity without that repair.
\end{proof}
